\section*{الفصل \ref{ch:b2:climate-variability} --- التقلّب المناخي}

\begin{solution}{exo:b2:climate-variability:1}
(أ) $29.27 - 26.71 = \qty{2.56}{\celsius}$ في 2026؛ و$28.79 - 26.73 = \qty{2.06}{\celsius}$ في 1997.
(ب) كان سبتمبر 2026 الشهر الأدفأ، بفارق \qty{0.48}{\celsius}، وصاحب الشذوذ الأكبر أيضًا، بفارق
\qty{0.50}{\celsius}. (ج) تُقارن كلّ سنة بفترة أساس طولها 30 سنة متمركزة عليها، تُزاح
كلّ خمس سنوات؛ \omterm{def:b2:climate-variability:modes}{والمعدّل المناخي} لشهر سبتمبر في عقدَي 1980 و1990 يختلف قليلًا عن نظيره في
عقدَي 2010 و2020.
\end{solution}

\begin{solution}{exo:b2:climate-variability:2}
(أ) JJA: $(0.79 + 0.95 + 1.28)/3 = 1.007$؛ JAS: $(0.95 + 1.28 + 1.55)/3 = 1.260$؛ ASO: $(1.28 +
1.55 + 1.68)/3 = 1.503$؛ SON: $(1.55 + 1.68 + 2.01)/3 = 1.747$؛ OND: $(1.68 + 2.01 + 2.02)/3 =
1.903$، وكلّها بوحدة \unit{\celsius}. (ب) مقرّبةً إلى منزلة عشرية واحدة: 1.0، 1.3، 1.5، 1.7، 1.9، وهي قيم
الجدول. (ج) في ASO 2023، عند 1.5.
\end{solution}

\begin{solution}{exo:b2:climate-variability:3}
(أ) يتتبّع المؤشّر انحراف الضغط في تاهيتي مطروحًا منه الضغط في داروين: فإذا كان دون المتوسط
في تاهيتي وفوقه في داروين، كان سالبًا. (ب) \omterm{def:b2:climate-variability:enso}{إل نينيو}. (ج) رياح تجارية شرقية أضعف، بل
رياح غربية قرب السطح في الغرب، ومطر أقلّ من المعتاد فوق إندونيسيا، لأن
الحمل العميق انتقل شرقًا.
\end{solution}

\begin{solution}{exo:b2:climate-variability:4}
(أ) $804.1/880.6 = \qty{91.3}{\percent}$ من المتوسط: ضمن المجال من 90 إلى \qty{110}{\percent}، فهو
موسم عادي، وإن كان في قسمه الأدنى. (ب) مجموع الموسم $556 + 659 + 427 + 252 =
\qty{1894}{mm}$؛ ويجلب يوليو $659/1894 = \qty{34.8}{\percent}$ منه، أي أكثر من الثلث.
\end{solution}

\begin{solution}{exo:b2:climate-variability:5}
(أ) $\tau = 1025 \times 3992 \times 50/16.1 = \qty{1.271e7}{s} = \qty{147.1}{\day}$ لعمق
\qty{50}{m}، وضعف ذلك، \qty{294.2}{\day}، لعمق \qty{100}{m}. (ب) $\mathrm{e}^{-30.44/147.1} =
0.813$ و$\mathrm{e}^{-30.44/294.2} = 0.902$. (ج) زمن التخامد متناسب مع $h$: فالطبقة
المختلطة الأضحل تختزن حرارة أقلّ للتدفّق نفسه وتنسى الشذوذ أسرع، فتتراجع
ذاكرة المحيط، وقابلية التنبّؤ التي تمنحها.
\end{solution}

\begin{solution}{exo:b2:climate-variability:6}
(أ) ثلاثة أشهر تساوي 91.3 يوم: $1.2 \times \mathrm{e}^{-91.3/147.1} = \qty{0.65}{K}$؛ وستّة أشهر
تساوي 182.6 يوم: $1.2 \times \mathrm{e}^{-182.6/147.1} = \qty{0.35}{K}$. (ب) يتناقص الشذوذ
ببطء، فهو تخمين لشذوذ الشهر المقبل أفضل من الصفر. لكن التأثير العشوائي يضيف
حالات شذوذ جديدة كلّ يوم، مجهولة وقت التنبّؤ؛ وكلّما تناقص الشذوذ الابتدائي،
أخذت حصّة متزايدة، فينجرف التنبّؤ نحو \omterm{def:b2:climate-variability:modes}{المعدّل المناخي}، حيث لا يعود
مفيدًا.
\end{solution}

\begin{solution}{exo:b2:climate-variability:7}
(أ) $160 \times 86\,400 = \qty{1.38e7}{J}$، أي \qty{13.8}{MJ} لكلّ متر مربّع في اليوم. (ب) تحمل
الطبقة $(1000 - 850) \times 100/9.81 = \qty{1529}{kg}$ من الهواء لكلّ متر مربّع، بسعة حرارية
قدرها $1529 \times 1004.5 = \qty{1.536e6}{J.m^{-2}.K^{-1}}$؛ فتسخن بمقدار $1.38 \times
10^7/1.536 \times 10^6 = \qty{9.0}{K}$ في اليوم. (ج) تشعّ الطبقة الحرارة بالأشعة تحت الحمراء، نحو
الفضاء ونحو الهواء الذي فوقها؛ والهواء الدافئ يصعد ويختلط بالطبقة التي فوقه، أو
تحمله الرياح بعيدًا، ومنها تدفّق الهواء الأبرد الذي يسحبه المنخفض الحراري نفسه.
\end{solution}

\begin{solution}{exo:b2:climate-variability:8}
(أ) $(-1.93 - 1.11 - 1.98)/3 = -1.673$، أي نحو $-1.67$. (ب) إنه أدنى شتاء في السجلّ.
(ج) \omterm{def:b2:climate-variability:nao}{تذبذب شمال الأطلسي} سالب بقوّة: منخفض آيسلندا ومرتفع الأزور ضعيفان، وتيّار نفّاث أضعف ومسار عواصف
أبعد جنوبًا من المعتاد؛ وشتاء بارد جافّ في شمال أوروبا، ورطب في جنوب أوروبا،
واندفاعات أكثر للهواء البارد في شرق الولايات المتحدة.
\end{solution}

\begin{solution}{exo:b2:climate-variability:9}
(أ) $9.3 - 0.4 = \qty{8.9}{mm}$ في اليوم، أكثر من 2.5: ضمن النطاق. (ب) في نصف الكرة الجنوبي
يكون الصيف ديسمبر--فبراير: $6.1 - 4.2 = \qty{1.9}{mm}$ في اليوم، أقلّ من 2.5: خارجه،
مع أن المطر هناك أغزر منه في معظم أماكن النطاق. (ج) $1894/122 = \qty{15.5}{mm}$ في
اليوم؛ ولا تخرج كيرالا من النطاق إلا إذا تجاوز معدّل شتائها $15.5 - 2.5 = \qty{13.0}{mm}$ في اليوم.
(د) \omterm{def:b2:climate-variability:monsoon}{الرياح الموسمية} انعكاس فصلي: فالمكان الرطب طوال السنة، مثل المحطّة
الثانية، لا رياح موسمية فيه، والمكان الجافّ الذي له موسم مطر قصير وشديد فيه رياح موسمية.
\end{solution}

\begin{solution}{exo:b2:climate-variability:10}
(أ) بالاشتقاق، $\ddot T = a\dot H = -abT$: معادلة مذبذب توافقي تردّده
الزاوي $\omega = \sqrt{ab}$، ودوره $2\pi/\sqrt{ab}$. (ب) $2\pi/\sqrt{ab} = 4$ يعطي $ab =
(2\pi/4)^2 = \qty{2.47}{yr^{-2}}$. (ج) مع $T = u\,\mathrm{e}^{-rt}$ و$H = v\,\mathrm{e}^{-rt}$،
$\dot T = (\dot u - ru)\,\mathrm{e}^{-rt} = (av - ru)\,\mathrm{e}^{-rt}$، فيكون $\dot u = av$،
وبالمثل $\dot v = -bu$: فالزوج $(u, v)$ يتذبذب دون تخامد، و$(T, H)$ هو ذلك التذبذب
مضروبًا في $\mathrm{e}^{-rt}$. وفي دور واحد تنخفض السعة بالعامل $\mathrm{e}^{-0.25 \times
4} = \mathrm{e}^{-1} = 0.37$. (د) تأثير خارجي: التذبذبات العشوائية للرياح التجارية، التي
تدفع النظام مرّة بعد مرّة، \omterm{def:b2:climate-variability:bjerknes}{والتغذية الراجعة لبيركنس}، التي تجعل التخامد أضعف أو
حتى سالبًا.
\end{solution}

\begin{solution}{exo:b2:climate-variability:11}
(أ) $A = R^2 \,\Delta\lambda\,(\sin\ang{5} - \sin(\ang{-5})) = (6.371 \times 10^6)^2 \times 0.8727
\times 0.1743 = \qty{6.17e12}{m^2}$. (ب) $Q = 1025 \times 3992 \times 100 \times 2 \times 6.17
\times 10^{12} = \qty{5.05e21}{J}$. (ج) ستّة أشهر تساوي $182.6 \times 86\,400 = \qty{1.578e7}{s}$؛
ولكلّ متر مربّع، $8.18 \times 10^8/1.578 \times 10^7 = \qty{51.9}{W/m^2}$، تدوم نصف
سنة. (د) من المحيط نفسه: فالرياح التجارية الأضعف تصعد بماء بارد أقلّ في الشرق، والماء الدافئ
في الغرب، المختزن فوق طبقة التدرج الحراري الأعمق، ينتشر شرقًا؛ وكثير من الاحترار
قد يأتي من حرارة موجودة في المحيط أصلًا، نُقلت لا اكتُسبت.
\end{solution}

\begin{solution}{exo:b2:climate-variability:12}
(أ) تعيد الحلقة $0.3g$، ثم $0.3g^2$، وهكذا: فالمجموع $0.3(1 + g + g^2 + \dots) =
0.3/(1 - g)$، وهي متسلسلة هندسية متقاربة إذا كان $g < 1$: \qty{0.6}{K} و\qty{1.5}{K}
و\qty{3.0}{K} للقيم $g = 0.5$ و0.8 و0.9. (ب) تتباعد المتسلسلة: فينمو الشذوذ بلا
حدّ؛ ويجب أن توقف النموَّ عملياتٌ تُغفلها الحلقة، هي التخامد بالتدفّق الحراري
السطحي وتفريغ مخزون الماء الدافئ الذي يصفه مذبذب إعادة الشحن. (ج)
كلّما كان سطح البحر أدفأ، أطلق احترار صغير بسهولة أكبر الحمل العميق والمطر،
اللذين يحرّك تسخينهما الرياح؛ وذلك هو الفصل الذي تكون فيه آثار ENSO أشدّ.
\end{solution}

\begin{solution}{pb:b2:climate-variability:1}
\textbf{1.} $(0.26 + 0.75 + 1.13)/3 = \qty{0.713}{\celsius}$.
\textbf{2.} MJJ $(0.75 + 1.13 + 1.52)/3 = 1.133$؛ JJA $(1.13 + 1.52 + 1.80)/3 = 1.483$؛ JAS $(1.52 +
1.80 + 2.06)/3 = 1.793$.
\textbf{3.} ASO $(1.80 + 2.06 + 2.26)/3 = 2.040$؛ SON $(2.06 + 2.26 + 2.42)/3 = 2.247$؛ OND $(2.26
+ 2.42 + 2.36)/3 = 2.347$؛ NDJ $(2.42 + 2.36 + 2.33)/3 = 2.370$.
\textbf{4.} مقرّبةً إلى منزلة عشرية واحدة، 0.7، 1.1، 1.5، 1.8، 2.0، 2.2، 2.3 و2.4: قيم الجدول.
\textbf{5.} الشهر المنفرد يحمل ضجيج الطقس؛ أمّا ENSO فيدوم أشهرًا كثيرة، ومتوسط ثلاثة أشهر
يزيل كثيرًا من الضجيج مع إبقاء الإشارة البطيئة.
\textbf{6.} في JJA 1997، عند 1.483، المقرَّبة إلى 1.5؛ وبلغ ذروته 2.4 في NDJ.
\textbf{7.} من ASO 2018 إلى MJJ 2019، ومن SON 2019 إلى FMA 2020.
\textbf{8.} عشرة فصول، وستّة.
\textbf{9.} في ASO 2020؛ ففي 2020 وحدها لا يكون عند $-0.5$ أو دونها إلا أربعة فصول، من ASO إلى NDJ، أي
أقلّ بفصل من القاعدة، ويكملها DJF 2021.
\textbf{10.} نوبة \omterm{def:b2:climate-variability:enso}{إل نينيو} واحدة، من SON 2018 إلى MAM 2019، سبعة فصول (ينخفض AMJ 2019 إلى 0.37)؛
وتبدأ \omterm{def:b2:climate-variability:enso}{لا نينيا} في AMJ 2020 ($-0.55$)، قبل أربعة فصول من الموضع الذي يضعها فيه المؤشّر ONI، أي ASO.
\textbf{11.} لا: فالمؤشّر النسبي لا يتجاوز 0.29 أبدًا في أواخر 2019. فقد كان المحيط المداري كلّه
حينئذ أدفأ من فترة أساسه بنحو نصف درجة؛ وكانت منطقة نينيو-3.4 دافئة
بالقيم المطلقة، لكنها بالكاد أدفأ من المنطقة المدارية حولها.
\textbf{12.} إحدى عشرة.
\textbf{13.} 8، 7، 10، 4، 5، 6، 12، 5، 9 و3 سنوات.
\textbf{14.} أقصرها 3 سنوات (من 2023 إلى 2026)، وأطولها 12 (من 1997 إلى 2009).
\textbf{15.} $69/10 = 6.9$ سنة بالطريقتين: فمجموع الفواصل المتتالية هو الزمن من
البداية الأولى إلى الأخيرة، $2026 - 1957 = 69$ سنة.
\textbf{16.} لقد دام بالفعل خمسة فصول، من MAM إلى JAS، عند \qty{0.5}{\celsius} أو فوقها،
وبلغ ذروة 2.2، فوق 1.5؛ وبدايته، MAM 2026، هي الفصل الأول من التتابع، والفصول اللاحقة
لا يمكنها إلا أن تمدّ التتابع، لا أن تنقل بدايته.
\textbf{17.} لا، ولا حاجة إلى ذلك: فالنوبات البالغ عددها 25 بكلّ شدّاتها تأتي كلّ ثلاث سنوات تقريبًا، كما تقول صفحة الأسئلة الشائعة؛
والقويّة منها مجموعة جزئية، فهي أندر، وتتراوح فواصلها بين 3 و12 سنة.
\textbf{18.} $\tau = 1025 \times 3992 \times 50/16.1 = \qty{1.27e7}{s} = \qty{147}{\day}$.
\textbf{19.} الطبقة المختلطة: $1025 \times 3992 \times 50 = \qty{2.05e8}{J.m^{-2}.K^{-1}}$؛ والعمود الهوائي:
$(10^5/9.81) \times 1004.5 = 10\,194 \times 1004.5 = \qty{1.02e7}{J.m^{-2}.K^{-1}}$، أي أقلّ
بعامل 20.
\textbf{20.} $1.024 \times 10^7/(1025 \times 3992) = \qty{2.5}{m}$ من ماء البحر.
\textbf{21.} $1.024 \times 10^7/16.1 = \qty{6.36e5}{s} = 7.4$~يوم.
\textbf{22.} شذوذ الهواء وحده يتخامد في نحو أسبوع، وشذوذ الطبقة المختلطة في
نحو خمسة أشهر، والطبقة الدافئة فوق طبقة التدرج الحراري أعمق من ذلك: فالمحيط وحده يحتفظ
بالحرارة مدّة تكفي لحمل الشذوذ من فصل إلى الذي يليه، بينما يقدّم الغلاف الجوي
الرياح التي تقرن بينهما.
\textbf{23.} منذ 1950، بدأت نوبات \omterm{def:b2:climate-variability:enso}{إل نينيو} القويّة في المتوسط كلّ \textbf{6.9 سنة}،
أي نحو سبع سنوات.
\end{solution}
